\documentclass{situationist-preprint}

\title{Lying to the Murderer at the Door}
\msid{ETH-04-a}
\msfamily{ETH-04}
\msdiscipline{Normative and Applied Ethics}
\mstype{Principal manuscript}
\mskeywords{Kant; lying; murderer at the door; categorical imperative; universalizability; maxims; individuation of maxims}

\begin{document}
\maketitle

\begin{abstract}
A murderer asks where your friend, whom he intends to kill, is hiding. Kant held that you may not lie to him, and his defenders have usually tried to show that the verdict follows from the formula of universal law. We argue that one may lie, and that the Kantian argument fails at the universalizability step itself. The test is applied to a maxim, and its verdict depends on how the maxim is described. Kant's argument tests the maxim \enquote{to lie when it serves my purpose} or \enquote{to lie to avert a harm}, which do fail the test. But those are not the maxim on which the person at the door acts. We state a criterion of individuation, on which a maxim must include every circumstance the agent treats as a reason for her act and nothing she does not, and show that the maxim so individuated passes the contradiction in conception test and the contradiction in the will test. The criterion blocks the gerrymandered maxims that make the test too permissive. It follows that the formula of universal law, correctly applied, permits the lie, and that the case is no objection to the formula, only to Kant's application of it.
\end{abstract}

\section{The problem}

In a short essay of 1797, replying to Benjamin Constant, Kant considered a case in which a would-be murderer asks whether our friend, whom he is pursuing, has taken refuge in our house \citep{Kant1797}. Constant had argued that the duty to tell the truth is owed only to those who have a right to the truth, and that the murderer has none \citep{Constant1797}. Kant replied that truthfulness in statements that cannot be avoided is a formal duty to everyone, however great the disadvantage to oneself or another, and that a lie, even a well-meant one, does wrong to humanity in general by making declarations as such unworthy of credence. He added that the liar is responsible for the consequences of the lie, however unforeseen, whereas the truthful person is not responsible for what the murderer does with the truth.

Few readers, Kantian or not, accept the verdict, and the case has become a standard test of Kant's ethics. The responses fall into three groups.
\begin{enumerate}[label=(\roman*)]
\item The verdict follows from the categorical imperative, and the categorical imperative is therefore mistaken. This is the common view outside Kantian ethics.
\item The verdict does not follow from the formula of universal law, correctly applied, but follows from something else, either from the formula of humanity or from Kant's doctrine of right. \citet{Korsgaard1986} argued that the maxim of lying to the murderer passes the universalizability test, because the murderer does not know that we know his purpose, but that the formula of humanity, taken as an ideal, forbids the lie; in nonideal conditions, she held, the formula of universal law may serve as a guide and permit it.
\item The verdict is a misapplication of Kant's own principles, which permit the lie \citep{Paton1954}, or it concerns a narrower class of statements, such as declarations with juridical force, than the case is usually taken to involve \citep{Wood2008}.
\end{enumerate}

We defend a version of (iii), and we locate the misapplication precisely. The universalizability test is a test of maxims, and its verdict depends on how the maxim under test is described. Kant's argument, and the arguments of those who accept its conclusion, test a maxim that is described too coarsely. Under a principled criterion of individuation, the maxim of the person at the door passes the test. The point that maxims can be described at many levels of generality is old; \citet{Anscombe1958} remarked that Kant's rule is useless without a stipulation about what counts as a relevant description of an action, and the difficulty has been discussed at length by \citet{ONeill1975} and \citet{Herman1993}. What has been missing is a criterion that permits the fine description in this case without permitting everything.

\section{Framework}

\begin{definition}[Maxim]\label{d:maxim}
A \emph{maxim} is a subjective principle of action of the form \enquote{in circumstances $C$, to do $A$, in order to bring about $E$}, where $C$ is a type of circumstance, $A$ a type of act and $E$ an end.
\end{definition}

\begin{definition}[Fineness]\label{d:fine}
A maxim $\langle C, A, E\rangle$ is \emph{finer} than $\langle C', A, E\rangle$ if $C$ entails $C'$ and not conversely. A maxim is \emph{too coarse} for an action if some feature of the circumstances that the agent treats as a reason for performing $A$ is omitted from $C$.
\end{definition}

\begin{principle}[Relevance]\label{p:relevance}
The maxim of an action, for the purpose of the universalizability test, is the maxim whose circumstance-type $C$ includes every feature of the situation that the agent treats as a reason for doing $A$ in pursuit of $E$, and no feature that she does not so treat.
\end{principle}

\begin{principle}[Generality]\label{p:generality}
The features in a maxim's circumstance-type are general: they are specified without reference to particular persons, places or times.
\end{principle}

\begin{definition}[Universalizability]\label{d:ful}
A maxim passes the \emph{contradiction in conception test} if a world in which everyone acts on it is conceivable and, in that world, acting on it remains a way of bringing about its end. It passes the \emph{contradiction in the will test} if an agent who wills the end can, without inconsistency, will that everyone act on it. A maxim is \emph{universalizable} if it passes both tests.
\end{definition}

Definition~\ref{d:ful} follows the practical reading of the contradiction in conception test defended by \citet{Korsgaard1985}: a maxim fails if, universalized, it would be self-defeating, that is, the act would no longer serve the end. This is the reading on which the test gives Kant's own verdict on the lying promise, and we adopt it because it is the reading most favourable to the Kantian prohibition. Principle~\ref{p:relevance} is the core of our argument. It ties the maxim to the agent's practical reasoning, which is where Kant locates it: a maxim is the principle on which the agent in fact acts, not a description chosen by the evaluator. Principle~\ref{p:generality} is Kant's requirement that a maxim be the kind of thing that could be a law, and it excludes maxims indexed to particulars.

\section{The argument}

The person at the door, whom we call the householder, faces a man who has made known, or let it be inferred, that he means to kill the friend; the man does not know that the householder has inferred this, and he asks a question whose truthful answer would enable the killing. The householder lies in order to prevent the killing. What she treats as reasons are these features: the questioner intends to kill someone; the answer sought would enable him to do so; he believes that his intention is concealed from her; and a false answer would prevent the killing. She does not treat as a reason that a lie would be convenient, nor that it would avert some harm or other.

\begin{lemma}\label{l:coarse}
The maxims \enquote{to make a false statement when it serves my purpose} and \enquote{to make a false statement when doing so would avert a harm} fail the contradiction in conception test.
\end{lemma}

\begin{proof}
In a world in which everyone makes false statements whenever doing so serves his purpose, no statement can be relied upon by a hearer who does not know the speaker's purpose. Since a false statement serves its maker's purpose only by being believed, the act of making it would not bring about the end, and the maxim is self-defeating. The second maxim is narrower but fails for the same reason: harms that a false statement might avert are ubiquitous, and a hearer could never exclude that the speaker took his statement to avert one. Declarations in general would lose credence, which is Kant's point about \enquote{humanity in general}.
\end{proof}

\begin{lemma}\label{l:fine}
The maxim $M^\ast$: \enquote{when someone who intends to kill another, and who believes his intention concealed from me, asks me a question whose true answer would enable the killing, to answer falsely, in order to prevent the killing} passes the contradiction in conception test.
\end{lemma}

\begin{proof}
Consider a world in which everyone acts on $M^\ast$. In that world, a false statement is made only to someone whose murderous intention is known to the speaker. A questioner who believes his intention concealed, as the antecedent requires, therefore has no reason to doubt the answer he receives: from his point of view, the practice of lying to known murderers does not apply to him. The false answer is believed, and it serves the end.

Suppose, against this, that murderers in the universalized world come to suspect that their intentions are always detected, and discount every answer. Then a murderer can no longer use questions to locate his victim, and the killing is prevented by that fact. The end of $M^\ast$ is to prevent the killing, not to produce a false belief, and it is secured in either case. There is therefore no world in which universal adoption of $M^\ast$ makes acting on it fail to serve its end.
\end{proof}

\begin{lemma}\label{l:will}
The maxim $M^\ast$ passes the contradiction in the will test.
\end{lemma}

\begin{proof}
An agent who wills that killings of the kind in question be prevented can will that everyone answer falsely to those who would use the truth to kill. Nothing she must will as a rational agent, such as that others assist her when she is in need or that she be able to rely on the word of others in general, is threatened by the universal adoption of $M^\ast$, whose antecedent is satisfied only by questions put in furtherance of a killing. On the contrary, as a potential victim she has reason to will it.
\end{proof}

\begin{theorem}\label{t:main}
The maxim on which the householder acts, individuated by Principles~\ref{p:relevance} and~\ref{p:generality}, is universalizable. The formula of universal law therefore does not forbid the lie.
\end{theorem}

\begin{proof}
By Principle~\ref{p:relevance}, the householder's maxim includes the features she treats as reasons, which are those listed above; by Principle~\ref{p:generality}, they are specified generally. The result is $M^\ast$. By Lemmas~\ref{l:fine} and~\ref{l:will}, $M^\ast$ passes both tests.
\end{proof}

\begin{corollary}\label{c:diagnosis}
The argument that the formula of universal law forbids the lie depends on testing a maxim that is too coarse in the sense of Definition~\ref{d:fine}.
\end{corollary}

\begin{proof}
The maxims of Lemma~\ref{l:coarse} omit the features the householder treats as reasons, and they are the only maxims for which the test returns a prohibition. By Theorem~\ref{t:main}, the maxim individuated by Principle~\ref{p:relevance} does not fail.
\end{proof}

Theorem~\ref{t:main} establishes permissibility relative to the formula of universal law. That the lie is permissible all things considered requires, in addition, that no other valid principle forbid it. We consider the two candidates in the objections below, the formula of humanity and the duty of right, and argue that neither does. Granted that, the lie is permissible; whether it is also required is a further question we do not address, though the duty to aid a friend in mortal danger makes it plausible that it is.

\section{Objections and replies}

\begin{objection}[Gerrymandering]\label{o:gerry}
If maxims may be individuated finely, any act passes the test. The maxim \enquote{to make a false promise to a creditor named Smith on a Tuesday in order to obtain money} is unlikely to be self-defeating when universalized, since few others will be in exactly that position.
\end{objection}

\begin{reply}
The gerrymandered maxim violates both principles of individuation. It violates Principle~\ref{p:generality}, since it refers to a particular person and a particular day. And it violates Principle~\ref{p:relevance}, since no one who makes a false promise treats the creditor's name or the day of the week as a reason for doing so; the reason is that money is needed and the promise will produce it. Restored to the features the agent treats as reasons, the maxim is the lying promise, which fails. Fine individuation is not arbitrary individuation. The criterion makes the description of the maxim answer to the agent's reasons, and an agent's reasons cannot be multiplied at will to evade the test, since a feature counts as a reason for her only if her decision is sensitive to it. The householder would not lie if she believed the questioner meant no harm; the false promiser would make his promise to Jones on a Wednesday.
\end{reply}

\begin{objection}[Concealed reasons]\label{o:conceal}
An agent may treat a morally objectionable feature as a reason, and Principle~\ref{p:relevance} then requires including it, which may let her maxim pass. A person who lies only to members of a despised group, and only when they cannot detect it, may have a universalizable maxim.
\end{objection}

\begin{reply}
The contradiction in conception test is not the whole of the formula of universal law. The maxim \enquote{to deceive members of group $G$ when they cannot detect it, in order to gain an advantage} fails the contradiction in the will test: an agent who wills it as a law wills that she be deceived whenever she belongs to a group that others single out, and no agent can rationally exclude that. More generally, the formula of universal law was never meant to identify all wrongs by the contradiction in conception test alone, and our argument does not require that it do so. It requires only that the maxim of the householder pass both tests, which it does.
\end{reply}

\begin{objection}[Humanity in general]\label{o:right}
Kant's argument in the 1797 essay is not an application of the universalizability test but an argument from the doctrine of right: every lie harms humanity in general by undermining the trust on which contracts and rights depend \citep{Kant1797}. Theorem~\ref{t:main} does not touch it.
\end{objection}

\begin{reply}
The argument from humanity in general has the same structure as Lemma~\ref{l:coarse}. A lie harms humanity in general only if it contributes to a practice in which declarations as such cease to be credible. The practice to which the householder's lie belongs is the practice of lying to concealed murderers about the whereabouts of their victims. That practice does not make declarations as such incredible, and it does not undermine contracts, which are not made in furtherance of a killing. The argument from right therefore fails at the same point as the argument from universal law: it treats the lie as an instance of the coarse type. A further consideration makes the same point from the side of right. A person who asks in order to kill is attempting to use the householder's speech as an instrument of a wrong, and on Kant's own view of right a hindrance to a hindrance of freedom is consistent with freedom \citep{Kant1797b}. The lie hinders his use of her words to hinder the freedom of another.
\end{reply}

\begin{objection}[Formula of humanity]\label{o:fh}
Even if the lie passes the formula of universal law, it treats the murderer merely as a means, since it manipulates his beliefs in a way to which he cannot consent \citep{Korsgaard1986}.
\end{objection}

\begin{reply}
The formula of humanity forbids treating a person \emph{merely} as a means, and ETH-21 argues that the moral weight of the formula lies in this word. On the account given there, one treats a person merely as a means when one uses him in a way he could reasonably reject, and a person cannot reasonably reject a treatment on a ground that depends on his pursuing a wrongful end. The murderer's only ground for objecting to the lie is that it frustrates his plan to kill. The householder uses his belief, but she treats him as an end: she does no more to him than is needed to protect her friend, and she would not, for example, harm him if a lie sufficed. The lie therefore treats him as a means but not merely as one.
\end{reply}

\begin{objection}[Responsibility for consequences]\label{o:conseq}
Kant also argued that the liar is answerable for whatever follows from the lie: if the friend has slipped out of the house and the murderer, sent elsewhere, meets him there, the liar is responsible for the death, whereas the truthful householder would not have been \citep{Kant1797}.
\end{objection}

\begin{reply}
The argument assumes that the truthful householder is not responsible for what the murderer does with her answer, and so assumes that telling the truth in this case is permissible, which is what is in question. If giving the murderer the information is wrong, the truthful householder is responsible for her contribution to the killing. The asymmetry Kant asserts is therefore not an independent premise. Nor does the unlucky outcome of the lie show that the lie was wrong; it shows at most that the householder, acting permissibly, contributed through bad luck to a bad outcome, as one may when acting on any maxim.
\end{reply}

\section{What the result does not establish}

The argument does not show that lying is generally permissible, nor that the duty of truthfulness is weak; the coarse maxims of Lemma~\ref{l:coarse} fail the test, and the lies they describe are wrong. It does not show that the formula of universal law is empty or that it can be made to permit anything: Principles~\ref{p:relevance} and~\ref{p:generality} are designed to prevent that, and Objection~\ref{o:gerry} shows them at work. It does not settle whether Kantian ethics is correct, nor whether its verdicts could be stated by a theory of another form. ETH-02 argues that every deontological theory without dilemmas has a consequentialist counterpart, so that a constraint against lying, with whatever exceptions it has, could be represented either way. The present question is first-order and is settled by the content of the constraint, which on the correct individuation of maxims does not extend to the householder.

\section{Conclusion}

One may lie to the murderer at the door. Kant's argument to the contrary fails at the universalizability step, because it tests a maxim that omits the features for which the householder acts. Individuated by the reasons the agent acts on and specified in general terms, her maxim passes both the contradiction in conception and the contradiction in the will tests, and the arguments from the duty of right and from the formula of humanity fail for related reasons. The case is an objection to Kant's application of the formula of universal law, not to the formula.

\begin{thebibliography}{99}\small
\bibitem[Anscombe(1958)]{Anscombe1958} Anscombe, G.~E.~M. (1958). Modern moral philosophy. \emph{Philosophy}, 33(124), 1--19.
\bibitem[Constant(1797)]{Constant1797} Constant, B. (1797). \emph{Des r\'eactions politiques}. Paris.
\bibitem[Herman(1993)]{Herman1993} Herman, B. (1993). \emph{The Practice of Moral Judgment}. Harvard University Press.
\bibitem[Kant(1797a)]{Kant1797} Kant, I. (1797a). On a supposed right to lie from philanthropy. In \emph{Practical Philosophy}, trans. M.~J. Gregor, Cambridge University Press, 1996. Ak.~8:423--430.
\bibitem[Kant(1797b)]{Kant1797b} Kant, I. (1797b). \emph{The Metaphysics of Morals}. In \emph{Practical Philosophy}, trans. M.~J. Gregor, Cambridge University Press, 1996. Doctrine of Right, Introduction, \S D.
\bibitem[Korsgaard(1985)]{Korsgaard1985} Korsgaard, C.~M. (1985). Kant's formula of universal law. \emph{Pacific Philosophical Quarterly}, 66(1--2), 24--47.
\bibitem[Korsgaard(1986)]{Korsgaard1986} Korsgaard, C.~M. (1986). The right to lie: Kant on dealing with evil. \emph{Philosophy \& Public Affairs}, 15(4), 325--349.
\bibitem[O'Neill(1975)]{ONeill1975} O'Neill, O. (1975). \emph{Acting on Principle: An Essay on Kantian Ethics}. Columbia University Press.
\bibitem[Paton(1954)]{Paton1954} Paton, H.~J. (1954). An alleged right to lie: a problem in Kantian ethics. \emph{Kant-Studien}, 45, 190--203.
\bibitem[Wood(2008)]{Wood2008} Wood, A.~W. (2008). \emph{Kantian Ethics}. Cambridge University Press.
\end{thebibliography}

\end{document}
